\documentclass{article}
\usepackage{graphicx} % Required for inserting images
\usepackage{amsmath}
\usepackage{ytableau}
\title{Solutions to Worksheet \#3}
\author{}
\date{August 4, 2026}

\begin{document}

\maketitle
\begin{enumerate}
\item \textbf{Solution.}
The semistandard tableaux of shape $(2,2)$ give
\[
\boxed{s_{2,2}=m_{2,2}+m_{2,1,1}+2m_{1,1,1,1}.}
\]
For example, the coefficient $2$ is the number of standard Young
tableaux of shape $(2,2)$.

The Jacobi--Trudi identity gives
\[
\boxed{s_{2,2}=h_2^2-h_3h_1=h_{2,2}-h_{3,1}.}
\]
Since $(2,2)$ is self-conjugate and the involution $\omega$ exchanges
$h_r$ with $e_r$, we also have
\[
\boxed{s_{2,2}=e_2^2-e_3e_1=e_{2,2}-e_{3,1}.}
\]
Finally, using the character values for the representation indexed by
$(2,2)$, or converting the preceding formulas with Newton's
identities, gives
\[
\boxed{s_{2,2}
=\frac{1}{12}p_{1,1,1,1}
+\frac14p_{2,2}
-\frac13p_{3,1}.}
\]

\item \textbf{Solution.}
\begin{enumerate}
\item Suppose a semistandard tableau of shape $\lambda$ and content
$\mu$ exists.  Because columns strictly increase, an entry in row
$r$ is at least $r$.  Therefore all entries among
$1,2,\ldots,r$ must lie in the first $r$ rows.  It follows that
\[
\mu_1+\cdots+\mu_r\leq\lambda_1+\cdots+\lambda_r
\qquad\text{for every }r.
\]
If $\mu>_{\mathrm{lex}}\lambda$, let $r$ be the first index at which
they differ.  Then $\mu_i=\lambda_i$ for $i<r$ and
$\mu_r>\lambda_r$, contradicting the displayed inequality.  Hence
\[
\boxed{\mu>_{\mathrm{lex}}\lambda\quad\Longrightarrow\quad
K_{\lambda,\mu}=0.}
\]

\item In a tableau of shape and content $\lambda$, all $\lambda_1$
copies of $1$ must lie in the first row, so that row consists entirely
of $1$'s.  After removing it, all $\lambda_2$ copies of $2$ must fill
the second row.  Continuing in this way shows that row $i$ must
consist entirely of $i$'s.  This tableau is semistandard and is the
only possibility.  Thus
\[
\boxed{K_{\lambda,\lambda}=1.}
\]
\end{enumerate}

\item \textbf{Solution.}
The monomial expansion of a Schur function is
\[
s_\lambda=\sum_{\mu\vdash k}K_{\lambda,\mu}m_\mu.
\]
Order the partitions of $k$ lexicographically.  Exercise 2 says that
this change-of-basis matrix is triangular, and its diagonal entries
are all $K_{\lambda,\lambda}=1$.  Its determinant is therefore $1$,
so it is invertible.  Since the monomial functions
$\{m_\mu:\mu\vdash k\}$ form a basis of $\Lambda^k$, the same is true
of
\[
\boxed{\{s_\lambda:\lambda\vdash k\}.}
\]

\item \textbf{Solution.}
Only the entries $1$ and $2$ are available.  Since columns strictly
increase, a tableau can have at most two rows.  Its shape must be
\[
(2k-r,r)\qquad\text{for some }0\leq r\leq k.
\]
In the $r$ columns of height two, the top entry must be $1$ and the
bottom entry must be $2$.  The remaining cells are all in the first
row; weak increase forces the remaining $k-r$ copies of $1$ to come
before the remaining $k-r$ copies of $2$.  Thus there is exactly one
tableau for each value of $r$.  The answer is
\[
\boxed{k+1.}
\]

\item \textbf{Solution.}
Let
\[
E(t)=\sum_{n\geq0}e_nt^n,\qquad
H(t)=\sum_{n\geq0}h_nt^n.
\]
The given identity is exactly the statement
\[
E(-t)H(t)=1.
\]
Suppose $\omega$ is defined by $\omega(e_n)=h_n$ and extended as an
algebra homomorphism.  Apply $\omega$ to the displayed identity:
\[
H(-t)\,\omega(H(t))=1.
\]
Replacing $t$ by $-t$ in the original identity also gives
$E(t)H(-t)=1$.  An inverse is unique, so
\[
\omega(H(t))=E(t).
\]
Comparing coefficients yields $\omega(h_n)=e_n$.  Thus $\omega^2$
fixes every $e_n$ (and every $h_n$).  Since the $e_n$ generate the
ring of symmetric functions,
\[
\boxed{\omega^2=\mathrm{id}.}
\]

\item \textbf{Solution.}
Since $h_r=s_r$, apply the Pieri rule to $s_3h_2$.  Adding a
horizontal strip of two boxes to the row $(3)$ produces exactly
$(5)$, $(4,1)$, and $(3,2)$.  Hence
\[
\boxed{h_{3,2}=s_5+s_{4,1}+s_{3,2}.}
\]

\item \textbf{Solution.}
Apply $\omega$ to the answer in Exercise 6.  The involution exchanges
$h$ and $e$ and sends $s_\lambda$ to $s_{\lambda'}$.  The conjugates
of $(5)$, $(4,1)$, and $(3,2)$ are respectively
$(1^5)$, $(2,1^3)$, and $(2,2,1)$.  Therefore
\[
\boxed{e_{3,2}=s_{1^5}+s_{2,1^3}+s_{2,2,1}.}
\]

\item \textbf{Solution.}
Newton's identities give
\[
p_3=3h_3-3h_2h_1+h_1^3.
\]
By the Pieri rule,
\[
h_3=s_3,\qquad
h_2h_1=s_3+s_{2,1},\qquad
h_1^3=s_3+2s_{2,1}+s_{1,1,1}.
\]
Substitution gives
\[
\boxed{p_3=s_3-s_{2,1}+s_{1,1,1}.}
\]

\item \textbf{Solution.}
The statement is false.  Let $\lambda=(2)$ and $\mu=(1,1)$.  There is
one semistandard tableau of shape $(2)$ and content $(1,1)$, namely a
row containing $1,2$.  Thus $K_{(2),(1,1)}=1$.  A column of shape
$(1,1)$ cannot be filled with two copies of $1$ because columns must
strictly increase, so $K_{(1,1),(2)}=0$.  Hence the two Kostka numbers
need not be equal.

\item \textbf{Solution.}
The partitions of $3$ give
\[
\boxed{\begin{aligned}
h_3&=s_3,\\
h_{2,1}&=s_3+s_{2,1},\\
h_{1,1,1}&=s_3+2s_{2,1}+s_{1,1,1}.
\end{aligned}}
\]
The pattern is Young's rule:
\[
\boxed{h_\lambda=\sum_{\mu\vdash3}K_{\mu,\lambda}s_\mu.}
\]
In general, the coefficient of $s_\mu$ in $h_\lambda$ is the number
of semistandard tableaux of shape $\mu$ and content $\lambda$.

\item \textbf{Solution.}
For $a\geq1$ and $b\geq0$, the Jacobi--Trudi determinant for the hook
$(a,1^b)$ is
\[
s_{a,1^b}=
\begin{vmatrix}
h_a&h_{a+1}&h_{a+2}&\cdots&h_{a+b}\\
1&h_1&h_2&\cdots&h_b\\
0&1&h_1&\cdots&h_{b-1}\\
\vdots&\ddots&\ddots&\ddots&\vdots\\
0&\cdots&0&1&h_1
\end{vmatrix}.
\]
Expand along the first row.  The minor corresponding to $h_{a+i}$ is
the Jacobi--Trudi determinant for $e_{b-i}$; the cofactor sign is
$(-1)^i$.  Therefore
\[
\boxed{s_{a,1^b}
=\sum_{i=0}^b(-1)^ih_{a+i}e_{b-i}.}
\]
This also proves the formula inductively: the case $b=0$ is
$s_a=h_a$, and increasing $b$ adds one more cofactor to the same
Jacobi--Trudi expansion.

\end{enumerate}
\end{document}
